\documentclass[lettersize,journal]{IEEEtran}
\usepackage{amsmath,amssymb}
\usepackage{array}
\usepackage{booktabs}
\usepackage{graphicx}
\usepackage{cite}
\usepackage{url}
\usepackage{xcolor}
\hyphenation{hetero-geneous accel-erator coher-ency}

\begin{document}

\title{AIRTOS: An Evidence-Constrained Runtime for Heterogeneous Edge AI Systems}

\author{Anonymous Authors}

\maketitle

\begin{abstract}
% To be written from idea.md and experimental_log.md.
\end{abstract}

\begin{IEEEkeywords}
Edge AI, heterogeneous computing, real-time systems, admission control, runtime assurance, fault recovery.
\end{IEEEkeywords}

\section{Introduction}
Edge inference increasingly crosses a boundary that conventional deployment
interfaces leave implicit. A compiler emits an execution plan, yet the plan
enters a live system whose provider health, contiguous memory, cache ownership,
outstanding device commands, and admitted workload can differ from the
conditions under which that plan was produced. Moreover, an inference is not
necessarily one schedulable thread. It can be a directed acyclic graph (DAG) of
CPU, vector, accelerator, and DMA segments with resource-specific execution
bounds and asynchronous completions. Consequently, structural validity and
numerical correctness alone do not establish that a plan is eligible to run in
the current state.

Real-time scheduling theory makes the assumptions behind timing guarantees
explicit. Classical earliest-deadline-first (EDF) results and sporadic-task
feasibility analysis depend on bounded execution and defined arrival models
\cite{liu1973scheduling,baruah1990preemptively}; work on execution-time
assurance further shows that the strength of a timing claim is inseparable from
the bounds supplied to the scheduler \cite{vestal2007preemptive}. Heterogeneous
task runtimes and schedulers already express DAG dependencies, locality, and
resource types \cite{topcuoglu2002performanceeffective,augonnet2011starpu,
bauer2012legion,rossbach2011ptask,lin2022typeaware}. Accelerator-sharing
systems likewise address preemption, spatial partitioning, memory pressure,
and multi-model quality of service
\cite{choi2020prema,ghodrati2020planaria,kim2023moca,kim2023dream}. These
mechanisms provide essential foundations, but they normally begin after the
workload and its execution bounds have been accepted as valid inputs.

This paper presents AIRTOS, an evidence-constrained runtime governance layer
for heterogeneous edge-AI plans. AIRTOS treats a compiler plan as a first-class
admission object and evaluates a joint predicate over parsing, binding,
applicability domain, evidence coverage, provider availability, memory,
coherency, schedulability, and recoverability. Admission probes an arena lease,
captures the current schedule, evaluates a conservative multi-resource
simulation, and commits memory and job state only if their generations remain
unchanged. During execution, dependency-ready segments enter stable
per-resource EDF queues. Plan-declared ownership transitions govern cache and
DMA operations, while a device epoch and dispatch cookie scope completion
events. Recovery consumes a bounded policy, quarantines an exhausted resource,
and subjects a fallback to the same evidence, provider, lease, and schedule
checks before execution.

The design follows the separation principle used in runtime assurance: a
complex producer may propose an action, but a smaller consumer-side mechanism
decides whether that action is admissible in the present state
\cite{seto1998simplex,hobbs2023runtime}. Runtime observations are similarly
restricted. Trace records can select a subsequent compiler experiment, but
they cannot elevate evidence or install a new plan directly; every candidate
must return through verification and admission. This preserves the distinction
between observing an execution and establishing a reusable claim, consistent
with runtime-verification practice \cite{leucker2009brief}.

We evaluate AIRTOS through four claim-oriented cores and a short physical
hardware-in-the-loop run. The frozen corpus covers package and admission
decisions, concurrent memory--schedule transactions, independent scheduling
oracle comparisons, arena leases, physical DMA transfers with negative
coherency controls, stale-event and recovery state transitions, fallback
gates, and trace classification. On one CanMV-K230-LP4 V3.0 board, one million
physical transfers produced no byte mismatch, and omission controls detected
all 800 intentionally missing ownership operations. The timing audit is
equally important: CPU and RVV maximum observations exceeded their registered
10 and 4~$\mu$s execution bounds, respectively. AIRTOS therefore identifies
the tested plan as outside its timing contract rather than converting favorable
percentiles into an unsupported hard-real-time claim.

This work makes four contributions. First, it defines evidence-constrained
admission that combines plan qualification with live resource, memory,
coherency, scheduling, and recovery state. Second, it implements an optimistic
atomic lease-and-schedule transaction for concurrent admission. Third, it
unifies dependency-aware resource governance, explicit ownership transitions,
and epoch-scoped completion and recovery semantics. Fourth, it provides a
finite-domain evaluation that reports both invariant checks and a negative
timing-contract result. The contribution is this governed composition and its
consumer-checkable boundaries; EDF, heterogeneous DAG execution, accelerator
sharing, and DMA synchronization remain established mechanisms.

\section{Background and Related Work}
\subsection{Real-Time and Heterogeneous DAG Scheduling}
Liu and Layland establish the classical uniprocessor foundation for EDF under
periodic-task assumptions \cite{liu1973scheduling}, while processor-demand
analysis extends feasibility reasoning to sporadic workloads
\cite{baruah1990preemptively}. Vestal makes execution-time assurance an
explicit scheduling dimension \cite{vestal2007preemptive}. AIRTOS retains this
discipline: its deadline statement is conditional on registered segment costs,
non-preemptive blocking, modeled arrivals, coherency charges, and recovery
charges. Its simulator rechecks the deadlines of all admitted jobs after a
candidate insertion instead of treating EDF priority as a timing proof.

HEFT schedules precedence-constrained tasks across heterogeneous processors
using rank-based list scheduling \cite{topcuoglu2002performanceeffective}.
StarPU combines task scheduling with heterogeneous data management
\cite{augonnet2011starpu}; Legion expresses independence and locality through
logical regions \cite{bauer2012legion}; and PTask makes accelerator work and
dataflow visible to the operating system \cite{rossbach2011ptask}. Type-aware
federated scheduling gives real-time semantics to typed heterogeneous DAGs
\cite{lin2022typeaware}. AIRTOS does not present DAG queues or per-resource EDF
as new scheduling algorithms. It places these mechanisms downstream of an
evidence and applicability gate and couples their admission result to a
generation-checked memory transaction and recovery state.

\subsection{Accelerator Runtimes and Multi-Tenant AI}
Accelerator sharing spans different hardware assumptions. PREMA exploits
preemptible NPU execution \cite{choi2020prema}, whereas Planaria partitions an
accelerator spatially among tenants \cite{ghodrati2020planaria}. MoCA adapts
execution around multi-tenant memory behavior \cite{kim2023moca}, and DREAM
targets dynamic real-time multi-model edge workloads \cite{kim2023dream}.
Salus exposes fine-grained GPU sharing primitives and distinguishes persistent
from transient memory \cite{yu2019salus}; Clockwork pursues predictable DNN
serving through controlled execution and explicit timing information
\cite{gujarati2020serving}. These systems establish that scheduling, memory,
and latency must be governed jointly, but their admitted unit is generally a
model or execution request rather than a plan carrying separately checkable
domain and evidence obligations. AIRTOS addresses that qualification boundary
and does not compare throughput across incompatible accelerator capabilities.

Embedded inference systems demonstrate how tightly runtime design is coupled
to constrained memory and platform support. TensorFlow Lite Micro provides an
embedded inference runtime for small devices \cite{david2021tensorflow}, and
MLPerf Tiny establishes reproducible measurement practices for such platforms
\cite{banbury2021mlperf}. AIRTOS focuses on a different layer: deciding whether
a particular evidence-carrying heterogeneous plan may enter the live runtime,
then maintaining its memory, timing, ownership, and recovery obligations.

\subsection{Runtime Assurance, Coherency, and Recovery}
Simplex separates a high-performance controller from a trusted baseline and
decision module \cite{seto1998simplex}; modern runtime-assurance work generalizes
that structure as a safety filter around complex components
\cite{hobbs2023runtime}. Runtime verification studies how execution events are
checked against specified properties \cite{leucker2009brief}. AIRTOS applies
the same separation at the compiler--runtime boundary: a producer supplies a
plan and evidence, while the runtime independently checks eligibility against
current state. Its trace is observational and can only propose a new
experiment, preventing an execution from certifying its own replacement.

DMA synchronization specifications already define fences and ownership
transitions for asynchronous devices \cite{documentation2026buffer,
documentation2026dma}. AIRTOS does not redefine those primitives. Instead, the
plan names the buffer range and required clean, barrier, and invalidate actions;
the runtime checks that the range lies within the active lease and invokes the
provider hooks at the ownership boundary. Completion attribution then adds an
orthogonal device-generation rule: an event is accepted only when resource,
epoch, active command, and cookie agree. This prevents an event associated with
an earlier device generation from completing a newly active segment.

Finally, the evaluation uses an independent scheduling oracle rather than
deriving expected answers from the runtime under test, following the broader
testing principle that oracle construction is part of the validity argument
\cite{barr2015oracle}. Together, these lines of work motivate AIRTOS as a
governance layer that composes established scheduling, task-runtime,
assurance, and DMA mechanisms around one consumer-checkable plan contract.

\section{System Model and Problem Formulation}
% To be filled by the section-writing stage.

\section{AIRTOS Design}
% To be filled by the section-writing stage.

\section{Implementation}
% To be filled by the section-writing stage.

\section{Experimental Methodology}
% To be filled by the section-writing stage.

\section{Evaluation}
% To be filled by the section-writing stage.

\section{Discussion}
% To be filled by the section-writing stage.

\section{Conclusion}
% To be filled by the section-writing stage.

\clearpage
\bibliographystyle{IEEEtran}
\bibliography{refs}

\end{document}
